\section{A moment theorem beyond the stipulated Gaussian grid}
\label{sec:general-moment}

I state the argument separately from the retained GBM and circuit
specialisation. Let $w_i>0$, $\sum_iw_i=1$, and let
$\ell_i(y)=a_i+b_i\cdot y$ be arbitrary affine log-price functions.
Set $\overline a=\sum_iw_i a_i$,
$\overline b=\sum_iw_i b_i$, $m_G=\overline a+\overline b\cdot Y$ and
$m_A=\log\sum_iw_i e^{\ell_i(Y)}$. The latent vector $Y$ can have any
joint law for which the moments specified below exist. Define its joint
moment-generating function by
$\Lambda(v)=\E e^{v\cdot Y}$ at the required arguments. Independence
is not assumed. Let $P:\R\to[0,1]$ be increasing and globally
$L$-Lipschitz, with $L<\infty$, and put
$X=P(m_A)-P(m_G)$.

\begin{theorem}\label{thm:general-moment}
Define $d_i=a_i-\overline a$ and $s_i=b_i-\overline b$. If
$\Lambda(s_i)$ and $\Lambda(s_i+s_k)$ are finite for every $i,k$, then
$0\leq X\leq1$ and
\begin{align}
 \E X^2&\leq\min\{1,L^2M\},\\
 M&=\sum_{i,k}w_iw_k e^{d_i+d_k}\Lambda(s_i+s_k)
       -2\sum_iw_i e^{d_i}\Lambda(s_i)+1.
\end{align}
In particular $M\geq0$. For an implemented output $\widehat X\in[0,1]$
with a proved pointwise error $|\widehat X-X|\leq\eta$,
\begin{equation}
 \operatorname{Var}(\widehat X)
 \leq\min\{1/4,(\sqrt{\min\{1,L^2M\}}+\eta)^2\}.
\end{equation}
Any separately proved mean-based variance bound can also enter this
minimum. The finite-law factorisation in the preceding section is the
special case $\Lambda(v)=\prod_jM_j(v_j)$.
\end{theorem}

\begin{proof}
For each $y$, $\sum_iw_i(d_i+s_i\cdot y)=0$. Jensen's inequality therefore
gives $T(y)=\sum_iw_i e^{d_i+s_i\cdot y}\geq1$.
Moreover $m_A-m_G=\log T\leq T-1$, as proved above. Monotonicity,
boundedness and the Lipschitz condition imply
$0\leq X\leq\min\{1,L(T-1)\}$. Squaring and taking expectations proves
the first assertion once $M=\E(T-1)^2$ is evaluated. The assumed
finite moment-generating functions make all sums integrable. Expanding
the finite sums for $T$ and $T^2$ gives the displayed expression for $M$,
which is nonnegative because it is the expectation of a square.
For the implemented output,
$\E\widehat X^2\leq\E(X+\eta)^2
\leq(\sqrt{\E X^2}+\eta)^2$, using
$\E X\leq\sqrt{\E X^2}$. Subtracting its nonnegative squared mean
bounds its variance by that second moment. Finally any variable in
$[0,1]$ has variance at most $1/4$, because
$\E \widehat X^2\leq\E\widehat X$ and
$p-p^2\leq1/4$ for $p\in[0,1]$.
\end{proof}

This result can certify conditional basket-payoff sampling under Gaussian
mixtures, finite scenario laws or other latent models with the stated joint
moments. A suitable kernel and joint moment certificate must be supplied
for each application. It does not justify treating an arbitrary financial
model as Gaussian, and it does not cover models with infinite required
exponential moments. The theorem is an original derivation in this
project; its theoretical novelty remains subject to comparison with
published moment and variance-reduction results.

\section{Weak dispersion and the variance lost in a canonical flag}
\label{sec:weak-dispersion}

I first state a distribution-general limit. The argument concerns
financial random variables before any fixed grid or arithmetic error.
It supplies an independent limit experiment, rather than extrapolating
the original 27 finite configurations.

\begin{theorem}\label{thm:general-weak-dispersion}
Let $V$ be an arbitrary $\R^p$-valued latent random vector, let
$w_i>0$ sum to one, and fix $b_i\in\R^p$ with
$\sum_iw_i b_i=0$. Write $B=\max_i\|b_i\|$ and
$\Sigma_b=\sum_iw_i b_i b_i^{\mathsf T}$.
For $0<|\varepsilon|\leq\varepsilon_0$, suppose deterministic coefficients
$a_i(\varepsilon)$ satisfy
\begin{equation}
 \sum_iw_i a_i(\varepsilon)=0,\qquad
 |a_i(\varepsilon)|\leq C\varepsilon^2
\end{equation}
for a finite common $C$. Define
\begin{equation}
 T_\varepsilon=\sum_iw_i
        e^{\varepsilon b_i\cdot V+a_i(\varepsilon)},\qquad
 J_\varepsilon=\log T_\varepsilon.
\end{equation}
Let the finite random base $m_\varepsilon$ converge almost surely to
a finite $m_0$, and let $P:\R\to[0,1]$ be increasing and continuously
differentiable, with $|P'|\leq L<\infty$. Assume
\begin{equation}\label{eq:general-dispersion-integrability}
 \E\{(1+\|V\|^4)
                 e^{2\varepsilon_0B\|V\|}\}<\infty.
\end{equation}
Then
\begin{equation}\label{eq:general-weak-L2}
 \frac{P(m_\varepsilon+J_\varepsilon)-P(m_\varepsilon)}
      {\varepsilon^2}
 \longrightarrow
 \tfrac12P'(m_0)V^{\mathsf T}\Sigma_b V=:H_0
 \quad\hbox{in }L^2.
\end{equation}
The mean is $\varepsilon^2\E H_0+o(\varepsilon^2)$ and the variance is
$\varepsilon^4\operatorname{Var}(H_0)+o(\varepsilon^4)$.
The respective $\Theta$ orders require
$\E H_0>0$ and $\operatorname{Var}(H_0)>0$.
\end{theorem}

\begin{proof}
Let $d_{i,\varepsilon}=\varepsilon b_i\cdot V+a_i(\varepsilon)$.
The assumptions imply $\sum_iw_i d_{i,\varepsilon}=0$ pointwise.
Jensen's inequality gives $T_\varepsilon\geq1$.
For fixed $V$, $d_{i,\varepsilon}/\varepsilon\to b_i\cdot V$, because
$|a_i(\varepsilon)/\varepsilon|\leq C|\varepsilon|\to0$.
Taylor's formula with integral remainder gives
\begin{equation}
 \frac{T_\varepsilon-1}{\varepsilon^2}
 =\sum_iw_i\frac{d_{i,\varepsilon}^2}{\varepsilon^2}
                 \int_0^1(1-t)e^{t d_{i,\varepsilon}}\,dt
 \longrightarrow\tfrac12V^{\mathsf T}\Sigma_b V.
\end{equation}
The centring hypothesis cancels the complete linear term, including
the order-$\varepsilon^2$ intercepts; no individual limit of
$a_i(\varepsilon)/\varepsilon^2$ is required.
Moreover
\begin{equation}
 0\leq\frac{T_\varepsilon-1}{\varepsilon^2}
 \leq\tfrac12(B\|V\|+C\varepsilon_0)^2
           e^{\varepsilon_0B\|V\|+C\varepsilon_0^2}
 =:D(V).
\end{equation}
The square of $D(V)$ is at most a finite constant times
$(1+\|V\|^4)e^{2\varepsilon_0B\|V\|}$ and is therefore integrable
by~\eqref{eq:general-dispersion-integrability}.
Since $0\leq\log T_\varepsilon\leq T_\varepsilon-1$ and
$\log(1+z)/z\to1$ at zero, $J_\varepsilon/\varepsilon^2$ has the
same pointwise limit and $L^2$ dominator.
The fundamental theorem of calculus writes the left-hand side
of~\eqref{eq:general-weak-L2} as
\begin{equation}
 \frac{J_\varepsilon}{\varepsilon^2}
       \int_0^1P'(m_\varepsilon+tJ_\varepsilon)\,dt.
\end{equation}
The integral converges almost surely to $P'(m_0)$ by continuity
and the base convergence, and its magnitude is at most $L$.
Dominated convergence with the integrable square $L^2D(V)^2$ proves
the $L^2$ limit. Cauchy--Schwarz proves convergence of first moments.
For the second moments I use
$|\E U^2-\E W^2|\leq\|U-W\|_2(\|U\|_2+\|W\|_2)$.
Subtracting the convergent squared means proves the variance expansion.
\end{proof}

Bounded latent scenario laws, centred Gaussian laws and suitable
exponentially integrable mixture laws can satisfy these hypotheses.
A moment-generating function finite at isolated arguments does not by
itself prove~\eqref{eq:general-dispersion-integrability}; the requisite
norm-exponential and polynomial moment must be checked. Degenerate
boundaries also matter. A symmetric one-dimensional two-point latent law
can make $V^{\mathsf T}\Sigma_bV$ constant. If its limiting base $m_0$
also has constant $P'(m_0)$, then $H_0$ is constant and the leading
fractional variance is zero. Constant quadratic form alone is
insufficient when $P'(m_0)$ remains random. The theorem does not replace
that boundary by a positive variance coefficient.

I next vary cross-sectional dispersion while retaining Gaussian
risk-neutral drifts. This is a financial specialisation of the general
theorem, and I retain its exact drift algebra and explicit constants.
Let $Y\sim N(0,I_p)$, fix positive weights, vectors
$b_i\in\R^p$ with $\sum_iw_i b_i=0$, a common residual loading
$R_0\in\R^p$, $\lambda>0$, and $c=r-d$. Set
\begin{align}
 R_i(\varepsilon)&=R_0+\varepsilon b_i,\qquad
 D_b=\sum_iw_i\|b_i\|^2,\\
 \mu_i(\varepsilon)&=c-\tfrac12\{
           \|R_i(\varepsilon)\|^2+\lambda^2\},\\
 \Sigma_b&=\sum_iw_i b_i b_i^{\mathsf T},\qquad
 Q(v)=v^{\mathsf T}\Sigma_b v.
\end{align}
The normalised increasing conditional payoff kernel $P:\R\to[0,1]$ is
continuously differentiable with globally bounded derivative.
The capped-lognormal kernel satisfies these conditions at every
$\lambda>0$. Define the arithmetic and geometric log-baskets
$m_{A,\varepsilon}$ and $m_{G,\varepsilon}$ as above, and
$X_\varepsilon=P(m_{A,\varepsilon})-P(m_{G,\varepsilon})$.

\begin{theorem}\label{thm:weak-dispersion}
As $\varepsilon\to0$ through nonzero real values,
\begin{equation}\label{eq:weak-L2}
 \frac{X_\varepsilon}{\varepsilon^2}
 \longrightarrow
 H_0(Y):=\tfrac12P'(m_0(Y))Q(Y-R_0)
 \quad\hbox{in }L^2,
\end{equation}
where
$m_0(Y)=c-\tfrac12(\|R_0\|^2+\lambda^2)+R_0\cdot Y$.
Consequently, with $p_\varepsilon=\E X_\varepsilon$,
\begin{align}
 p_\varepsilon&=\varepsilon^2\E H_0+o(\varepsilon^2),\\
 \operatorname{Var}(X_\varepsilon)
   &=\varepsilon^4\operatorname{Var}(H_0)+o(\varepsilon^4).
\end{align}
If $\E H_0>0$ and $\operatorname{Var}(H_0)>0$, these orders are
respectively $\Theta(\varepsilon^2)$ and $\Theta(\varepsilon^4)$.
\end{theorem}

\begin{proof}
Weighted centring of the loading and the drift gives the exact identities
\begin{align}
 m_{G,\varepsilon}
 &=m_0(Y)-\tfrac12\varepsilon^2D_b,\\
 d_{i,\varepsilon}(Y)
 &:=\mu_i(\varepsilon)-\sum_kw_k\mu_k(\varepsilon)
              +(R_i(\varepsilon)-R_0)\cdot Y\\
 &=\varepsilon u_i(Y)+\varepsilon^2v_i,\\
 u_i(Y)&=b_i\cdot(Y-R_0),\qquad
 v_i=-\tfrac12(\|b_i\|^2-D_b).
\end{align}
Both $\sum_iw_i u_i$ and $\sum_iw_i v_i$ are zero. Thus
$T_\varepsilon=\sum_iw_i e^{d_{i,\varepsilon}}\geq1$ and
$J_\varepsilon=m_{A,\varepsilon}-m_{G,\varepsilon}
=\log T_\varepsilon\geq0$.
At each fixed $Y$, Taylor's formula gives
\begin{equation}
 \frac{T_\varepsilon-1}{\varepsilon^2}
 \longrightarrow \tfrac12\sum_iw_i u_i(Y)^2
 =\tfrac12 Q(Y-R_0).
\end{equation}
For a detailed domination argument, the integral Taylor remainder
$e^z-1-z=z^2\int_0^1(1-t)e^{tz}\,dt$ implies
$0\leq e^z-1-z\leq z^2e^{|z|}/2$ for real $z$.
Since the weighted linear terms vanish,
\begin{equation}
 0\leq\frac{T_\varepsilon-1}{\varepsilon^2}
 \leq \tfrac12\sum_iw_i
       (|u_i|+|v_i|)^2e^{|u_i|+|v_i|}
 =:D(Y),\qquad0<|\varepsilon|\leq1.
\end{equation}
Each $u_i$ is affine in $Y$. Hence
$D(Y)\leq C(1+\|Y\|^2)e^{C(1+\|Y\|)}$ for a finite constant $C$
depending only on the fixed coefficients. Its square is integrable under
the Gaussian law: completing the square in the radial Gaussian exponent
dominates any fixed linear exponential times a polynomial by an integrable
Gaussian with a smaller negative quadratic coefficient.
Since $0\leq\log T_\varepsilon\leq T_\varepsilon-1$ and
$\log(1+z)/z\to1$ as $z\downarrow0$, the same pointwise limit and
$L^2$ domination hold for $J_\varepsilon/\varepsilon^2$.

The fundamental theorem of calculus gives
\begin{equation}
 \frac{X_\varepsilon}{\varepsilon^2}
 =\frac{J_\varepsilon}{\varepsilon^2}
   \int_0^1P'(m_{G,\varepsilon}+tJ_\varepsilon)\,dt.
\end{equation}
At each $Y$, the integral converges to $P'(m_0(Y))$ by continuity.
Its absolute value is at most the global derivative bound $L$.
The product is dominated in square by $L^2D(Y)^2$.
Dominated convergence proves~\eqref{eq:weak-L2}.
An $L^2$ limit gives convergence of first moments by
Cauchy--Schwarz and convergence of second moments by
$|\E U^2-\E V^2|\leq\|U-V\|_2(\|U\|_2+\|V\|_2)$.
Subtracting the convergent squared means proves the variance expansion.
The asserted $\Theta$ orders follow when their leading coefficients are
strictly positive.
\end{proof}

\begin{corollary}\label{cor:weak-centred}
If $R_0=0$, $\Sigma_b\ne0$ and
$a=P'(c-\lambda^2/2)>0$, then
\begin{equation}
 \E H_0=\tfrac a2\operatorname{tr}\Sigma_b,\qquad
 \operatorname{Var}(H_0)=\tfrac{a^2}{2}
                         \operatorname{tr}(\Sigma_b^2)>0.
\end{equation}
\end{corollary}

\begin{proof}
Here $m_0$ is constant and $H_0=aY^{\mathsf T}\Sigma_bY/2$.
The real symmetric positive semidefinite matrix $\Sigma_b$ has an
orthogonal diagonalisation with eigenvalues $\rho_j\geq0$.
The Gaussian density is invariant under an orthogonal transformation,
so $Y^{\mathsf T}\Sigma_bY$ has the law of
$\sum_j\rho_j Z_j^2$ with independent standard normals $Z_j$.
Integration by parts using $\varphi'(z)=-z\varphi(z)$ gives
$\E Z_j^2=1$ and $\E Z_j^4=3$; all boundary terms vanish because
a polynomial times the Gaussian density vanishes at infinity.
Thus $\operatorname{Var}(Z_j^2)=2$. Independence now gives
$\E Q=\sum_j\rho_j$ and
$\operatorname{Var}(Q)=2\sum_j\rho_j^2$, which proves the formulas.
At least one $\rho_j$ is positive when $\Sigma_b\ne0$.
\end{proof}

Let a canonical flag be
$B_\varepsilon=\mathbf1\{U<X_\varepsilon\}$ with an independent
uniform $U$ on $(0,1)$. The flag has exactly the Bernoulli law with
parameter $p_\varepsilon$. Conditioning on $Y$ gives the exact identity
\begin{equation}\label{eq:flag-variance-loss}
 \operatorname{Var}(B_\varepsilon)
 =\operatorname{Var}(X_\varepsilon)+
                        \E[X_\varepsilon(1-X_\varepsilon)].
\end{equation}
It follows either by expanding $\E B_\varepsilon^2=p_\varepsilon$,
or by the conditional-variance formula since
$\E(B_\varepsilon\mid Y)=X_\varepsilon$ and
$\operatorname{Var}(B_\varepsilon\mid Y)
=X_\varepsilon(1-X_\varepsilon)$.
Under the nondegeneracy conditions above,
$\operatorname{Var}(B_\varepsilon)
 =\varepsilon^2\E H_0+o(\varepsilon^2)$.
The fractional output has variance of order $\varepsilon^4$, whereas
the canonical single-shot flag has variance of order $\varepsilon^2$.

The implication concerns a defined estimation experiment. For the
arithmetic sample mean of $N$ independent draws,
the relative mean-square errors are exactly
$\operatorname{Var}(X_\varepsilon)/(Np_\varepsilon^2)$ and
$p_\varepsilon(1-p_\varepsilon)/(Np_\varepsilon^2)$.
At a fixed relative error target these scale respectively as a constant
divided by $N$, and as a positive constant divided by
$N\varepsilon^2$. Chebyshev's inequality turns a proved variance bound
and a proved positive lower bound for $p_\varepsilon$ into a sufficient
relative-error confidence count. The theorem is not a lower bound for
all quantum estimators, and an amplitude-estimation algorithm exploits
coherent powers rather than independent unamplified flag averages.
Nor do I claim that the present bounded-range Bernstein formula yields
a dispersion-independent relative-error count: its linear range term
can dominate. At a fixed absolute monetary tolerance, the residual
eventually becomes smaller than that tolerance, so omitting its
estimation can become valid. I therefore retain the relative-residual
experiment when interpreting these weak-dispersion orders.

This identifies a financial loss of variance information in a particular
canonical encoding, rather than a universal quantum disadvantage.
At fixed numerical precision a grid or arithmetic error floor can dominate
the limit. Transferring the continuous result to a sequence of implemented
outputs requires a separately proved total $L^2$ implementation and
latent-law error $o(\varepsilon^2)$. A pointwise arithmetic error
$\eta_\varepsilon=o(\varepsilon^2)$ supplies its arithmetic component,
but a mean-bias bound alone does not supply the required latent-law
$L^2$ component. The original finite confidence certificates remain
separate from this asymptotic statement.
The mechanism is a project-specific derivation from established
Rao--Blackwell conditioning and small-dispersion analysis. A defensible
claim of a new contribution requires the literature comparison and
finite-parameter error evidence retained elsewhere in this project.
